\section{NVMe over Fabrics：把本地队列延伸到远端控制器}

本地 NVMe 把 SQ、CQ 和数据缓冲放在主机 DRAM，控制器通过 PCIe DMA 访问这些结构。\term{NVMe over Fabrics}{NVMe-oF}保留“命令、队列对、CID、完成状态和命名空间”这套上层模型，却把传输边界改成 RDMA、TCP 或 Fibre Channel 等 fabric。主机不再允许远端控制器任意读取本机物理内存；传输层必须把命令 capsule、数据和 completion 按各自协议搬过网络。\cite{nvme-of}

\topic{连接建立先于 I/O 队列所有权}

主机先发现一个 NVMe subsystem 的传输地址，建立 fabric 连接，并通过 Connect 命令创建管理队列。完成控制器识别、属性读取和命名空间发现后，主机才建立一个或多个 I/O queue pair。每条连接都有队列身份、深度、超时和控制器关联；网络 socket 或 RDMA queue pair 建立成功，只表示传输通道可用，不表示 NVMe 控制器已经接受对应 I/O 队列。

NVMe/TCP 把命令、响应与数据组织为带长度和校验边界的 PDU。写操作可以由主机发送命令后继续发送数据，也可以由控制器先请求特定数据范围；读操作则由控制器返回数据，再给出命令完成。RDMA 传输可以使用已注册内存，让远端控制器通过受权限约束的 RDMA Read/Write 搬运数据。注册键相当于临时访问能力：命令结束、连接失效或缓冲区回收时必须撤销，不能把一次映射永久留给远端设备。

\begin{pdfdiagram}[x=1cm,y=1cm]
  % 三层水平布局分别表示主机、fabric 和远端子系统。
  \node[hw state, minimum width=2.55cm, align=center] (req) at (0.7,4.0) {块请求\\tag 与缓冲区};
  \node[hw control, minimum width=2.75cm, align=center] (hostq) at (4.0,4.0) {NVMe 主机\\I/O queue pair};
  \draw[hw bus] (req) -- (hostq);

  \node[hw block, minimum width=3.0cm, minimum height=1.9cm, align=center] (fabric) at (7.7,2.65)
    {RDMA / TCP fabric\\命令 capsule\\数据与完成 PDU};
  \draw[hw bus] (hostq.south east) -- (6.0,4.0) -- (fabric.north west);

  \node[hw control, minimum width=2.75cm, align=center] (ctrl) at (11.35,4.0) {远端控制器\\CID 与命名空间};
  \node[hw memory, minimum width=2.75cm, align=center] (media) at (11.35,1.25) {SSD / 阵列介质\\缓存与持久化域};
  \draw[hw bus] (fabric.north east) -- (9.45,4.0) -- (ctrl.south west);
  \draw[hw bus] (ctrl) -- (media);

  \node[hw state, minimum width=2.55cm, align=center] (done) at (0.7,1.25) {主机完成\\状态与残余长度};
  \draw[hw return] (media.west) -- (7.7,1.25) -- (fabric.south);
  \draw[hw return] (fabric.west) -- (4.0,1.25) -- (done);
\end{pdfdiagram}
\diagramnote{一笔 NVMe-oF I/O 的三个所有权域。主机队列发布请求，fabric 只负责按 RDMA/TCP 契约传递命令和数据，远端控制器才解释命名空间并访问介质。网络 ACK、RDMA completion 和 NVMe CQ completion 属于不同层级；只有最后一项能结束 NVMe 命令，持久化还要满足 FUA/Flush 语义。}

\topic{传输完成、命令完成与持久化完成必须分开}

TCP 确认表示字节已经由对端协议栈接收，RDMA 本地 completion 表示工作请求达到该传输操作的完成条件；两者都不能证明远端 NVMe 命令已经成功。控制器可能在接收全部写数据后报告介质错误，也可能把数据放入易失缓存后返回普通成功。主机只有收到带匹配 CID 的 NVMe completion，才能结束命令；若上层要求掉电后可见，还必须使用设备声明支持的 FUA 或 Flush，并把远端阵列内部复制、缓存和电源保护纳入持久化域。

超时最困难的情况是“结果未知”：主机没有收到 completion，但远端控制器可能已经执行写入。对普通覆盖写，重试通常仍需依赖上层幂等性和顺序日志；对带外部副作用的管理命令，不能把超时直接解释为失败。安全恢复要保留请求身份和连接 generation，先确定旧连接是否仍可能返回数据，再在新连接上查询状态或按协议重放。

\section{ANA、多路径与无双主切换}

同一 namespace 可以通过多个控制器和多条 fabric 路径到达。多路径的价值不只是提高带宽，更重要的是在链路、端口或控制器故障时维持单一 I/O 语义。主机要区分“TCP 连接还在”“控制器仍响应 Keep Alive”“该 namespace 经这条路径仍可访问”三个状态。

\term{非对称命名空间访问}{Asymmetric Namespace Access, ANA}给每个路径组报告访问状态。Optimized 表示当前路径适合正常 I/O；Non-Optimized 仍可访问，但可能多经过一次内部转发；Inaccessible 暂时不能访问；Persistent Loss 表示路径不会自行恢复；Change 表示状态正在迁移。状态改变事件只触发重新读取 ANA 日志，不能直接替主机完成所有未决请求。

\begin{longtable}{L{0.18\linewidth}L{0.24\linewidth}L{0.29\linewidth}L{0.19\linewidth}}
\toprule
观测事件 & 能确认的事实 & 不能立即确认 & 主机下一步 \\
\midrule
Keep Alive 成功 & 控制器管理路径仍能响应 & 某 namespace 的数据路径最优或可用 & 继续结合 ANA 与 I/O 完成判断 \\\tablerowrule
ANA Non-Optimized & namespace 仍可经该路径访问 & 性能、内部转发距离与持久化域相同 & 可继续或迁往 Optimized 路径 \\\tablerowrule
连接超时 & 当前传输路径无进展 & 未完成写一定未执行 & 冻结该路径，隔离迟到完成并选择替代路径 \\\tablerowrule
ANA Persistent Loss & 该路径不会按普通状态迁移自动恢复 & 其他控制器也不可访问 namespace & 撤销该路径资源，检查其余路径 \\\tablerowrule
新路径 I/O 成功 & 新控制器可完成新请求 & 旧路径没有迟到数据或完成 & 依 generation 拒绝旧事件并核对写顺序 \\
\bottomrule
\end{longtable}

\topic{一次安全故障切换的状态链}

假设路径 A 上有 64 笔 outstanding I/O，随后 Keep Alive 和 I/O 同时超时。主机先停止向 A 提交新命令，把 A 的 queue pair 标为 quiescing，并保存每笔请求的 CID、写入范围和持久化要求。随后撤销或失效 A 使用的传输资源，确保旧 RDMA key、socket 和 completion 不会被新队列复用。主机读取路径 B 的 ANA 状态；只有 B 对目标 namespace 可访问，才为 B 建立新 generation 的队列。

未完成的读可以在新缓冲区上重发，因为旧路径仍可能迟到写入原缓冲区。未完成的写则要按上层日志和幂等规则分类：从未发布的数据可以直接重发；执行结果未知的写要保留原顺序与事务身份；已经收到 NVMe 成功但尚未满足持久化屏障的写仍不能宣称落盘。新路径返回成功后，主机还要拒绝路径 A 的迟到 completion，并验证旧写不会覆盖已经转作他用的内存。

\begin{pdfdiagram}[x=1cm,y=1cm]
  \node[hw state, minimum width=2.6cm, align=center] (run) at (0.7,3.0) {路径 A 运行\\generation 12};
  \node[hw control, minimum width=2.6cm, align=center] (freeze) at (3.85,3.0) {超时：冻结提交\\保存 64 笔身份};
  \node[hw control, minimum width=2.6cm, align=center] (isolate) at (7.0,3.0) {隔离 A\\撤销 key / 队列};
  \node[hw state, minimum width=2.6cm, align=center] (ana) at (10.15,3.0) {读取 B 的 ANA\\确认可访问};
  \node[hw state, minimum width=2.6cm, align=center] (resume) at (13.3,3.0) {路径 B 重建\\generation 13};
  \draw[hw bus] (run) -- (freeze);
  \draw[hw bus] (freeze) -- (isolate);
  \draw[hw bus] (isolate) -- (ana);
  \draw[hw bus] (ana) -- (resume);

  \node[hw label, text width=12.2cm, align=center] at (7.0,1.15)
    {恢复验收：读使用新缓冲区；结果未知的写按日志重放；B 的 completion 只匹配 generation 13；A 的迟到数据与完成均被隔离；最后核对 Flush/FUA 所定义的持久化顺序。};
\end{pdfdiagram}
\diagramnote{多路径切换不是把指针从 A 改到 B，而是一次所有权迁移。旧传输资源先失效，新路径状态随后确认，最后才用新 generation 恢复 I/O；这样才能避免双主、CID 复用和迟到 DMA 污染。}

NVMe-oF 把设备放远了，但没有改变本书反复使用的判断方法：传输层负责把消息送到，NVMe 层负责把完成归属于命令，上层负责决定结果能否重试，持久化协议负责决定掉电后是否可见。任何“网络通了所以存储完成”的推断都跨越了尚未验证的边界。
